% Dosłowne opisy z arkusza wyniki/Analiza_porownawcza_ODER_23_2026.xlsx
% (numeracja laboratoriów: 1 -- QSI, 2 -- Eurofins, 3 -- Intertek, 4 -- Puławy)
{\small
\begin{longtable}{@{}lc>{\raggedright\arraybackslash}p{5.2cm}>{\raggedright\arraybackslash}p{5.6cm}@{}}
\caption{Dosłowne opisy zapachu i~smaku zgłoszone przez cztery laboratoria
dla tych samych partii miodu (sierpień 2026). Numeracja
laboratoriów (1--4) jest umowna i~zgodna z~tabelą~\ref{tab:panel_lab}.\label{tab:interlab_opisy}}\\
\toprule
\textbf{Próbka} & \textbf{Lab.} & \textbf{Zapach} & \textbf{Smak} \\
\midrule
\endfirsthead
\multicolumn{4}{c}{\tablename\ \thetable{} (cd.)}\\
\toprule
\textbf{Próbka} & \textbf{Lab.} & \textbf{Zapach} & \textbf{Smak} \\
\midrule
\endhead
\bottomrule
\endfoot
P1
 & 1 & nietypowy dla miodu & nietypowy dla miodu \\
 & 2 & przypominający mocz zwierzęcy, lekko typowy dla pożytku & słodki, kwiatowy, łagodny, lekko typowy dla pożytku \\
 & 3 & typowy dla miodu & słodki, owocowy \\
 & 4 & słaby, lekko kwaśny, charakterystyczny dla oznak fermentacji & słodki, łagodny \\
\midrule
P2
 & 1 & nietypowy dla miodu & nietypowy dla miodu \\
 & 2 & przypominający mocz zwierzęcy, lekko typowy dla pożytku & słodki, kwiatowy, łagodny, lekko metaliczny, lekko typowy dla pożytku \\
 & 3 & lekko typowy dla pożytku & słodki, owocowy, lekko typowy dla pożytku \\
 & 4 & nienaturalny, lekko kwiatowy & bardzo słodki, szczypiący \\
\midrule
P3
 & 1 & typowy dla pożytku & typowy dla pożytku \\
 & 2 & lekko typowy dla pożytku & słodki, kwiatowy, ziołowy, lekko typowy dla pożytku \\
 & 3 & typowy dla miodu & słodki, kwiatowy, owocowy, ziołowy \\
 & 4 & silny, podobny do wosku pszczelego & słodki, nieco ostry \\
\midrule
P4
 & 1 & słabo typowy dla pożytku & słabo typowy dla pożytku \\
 & 2 & lekko typowy dla miodu & słodki, owocowy, medyczny (apteczny), lekko metaliczny, lekko typowy dla miodu \\
 & 3 & lekko typowy dla pożytku & słodki, owocowy, miętowy, lekko typowy dla pożytku \\
 & 4 & słabo wyczuwalny, kwiatów rzepaku & słodki \\
\midrule
P5
 & 1 & słabo typowy dla pożytku & słabo typowy dla pożytku \\
 & 2 & lekko typowy dla pożytku & słodki, kwiatowy, medyczny (apteczny), lekko gorzki, miętowy, lekko typowy dla pożytku \\
 & 3 & typowy dla miodu & słodki, owocowy, miętowy \\
 & 4 & wyraźny, charakterystyczny dla oznak fermentacji & słodki, łagodny \\
\midrule
P6
 & 1 & nietypowy dla miodu & nietypowy dla miodu \\
 & 2 & brak & słodki, bez wyraźnego aromatu kwiatowego, nietypowy dla miodu \\
 & 3 & nietypowy dla miodu & słodki, owocowy, słodowy, nietypowy dla miodu \\
 & 4 & silny, charakterystyczny dla oznak fermentacji & słodki, przypominający palony karmel \\
\midrule
P7
 & 1 & nietypowy dla miodu & nietypowy dla miodu \\
 & 2 & nietypowy dla miodu & słodki, bez wyraźnego aromatu kwiatowego, lekko metaliczny, nietypowy dla miodu \\
 & 3 & nietypowy dla miodu & słodki, owocowy, słodowy, nietypowy dla miodu \\
 & 4 & silny, landrynkowy & bardzo słodki, nieco szczypiący \\
\end{longtable}
}
